(1) 在长度 \(dz\) 上用电压降与流向地电流分别得到 \(\partial_z V=-(R'+i\omega L')I\) 和 \(\partial_z I=-(G'+i\omega C')V\)。消去任一相量有 \(\partial_z^2V=\gamma^2V,\ \partial_z^2I=\gamma^2I\)，其中 \(\gamma=\sqrt{(R'+i\omega L')(G'+i\omega C')}\)，特性阻抗 \(Z_0=\sqrt{(R'+i\omega L')/(G'+i\omega C')}\)。选取被动线的物理平方根使 \(\Re\gamma\ge0\)，并与 \(\Re Z_0>0\) 的能流分支相匹配；右行波为 \(e^{-\gamma z}\)，虽然在 \(z\le0\) 上从左向右传播。

(2) 写 \(V(z)=V_+e^{-\gamma z}+V_-e^{\gamma z}\)，对应 \(I(z)=[V_+e^{-\gamma z}-V_-e^{\gamma z}]/Z_0\)。在负载 \(z=0\) 有 \(V(0)=Z_LI(0)\)，求得 \(\Gamma_L=V_-/V_+=(Z_L-Z_0)/(Z_L+Z_0)\)。反射为零当且仅当 \(Z_L=Z_0\)；一般有损线的 \(Z_0\) 是复数，仅让负载电阻等于 \(|Z_0|\) 并不保证无反射。有限传播距离上的反射振幅还需计入往返衰减因子。

(3) 无损线 \(R'=G'=0\) 时 \(\gamma=i\beta,\ \beta=\omega\sqrt{L'C'}\)，\(Z_0=\sqrt{L'/C'}\) 为正实数。两行波的时间平均净功率 \(P=(|V_+|^2-|V_-|^2)/(2Z_0)=P_{\rm in}(1-|\Gamma_L|^2)\)，而负载 \(P_L=\frac12\Re[V(0)I(0)^*]\) 等于它。开路 \(Z_L\to\infty\) 有 \(\Gamma=+1\)（电流相消）；短路 \(Z_L=0\) 有 \(\Gamma=-1\)（电压相消）；\(Z_L=Z_0\) 时 \(\Gamma=0\) 且全部入射功率交付负载。

(4) 令 \(Z=R'+i\omega L'\)，\(Y=G'+i\omega C'\)，乘积写成 \(-\omega^2L'C'[1-i(R'/\omega L'+G'/\omega C')+O(\epsilon^2)]\)。取平方根，得 \(\gamma\simeq\alpha+i\beta\)，其中 \(\beta\simeq\omega\sqrt{L'C'}\)，\(\alpha\simeq\frac{R'}2\sqrt{C'/L'}+\frac{G'}2\sqrt{L'/C'}\)。由于 \(\alpha\ge0\)，右行波振幅随传播距离衰减；若线路接近截止、\(\omega\) 太小或 \(R'/(\omega L')\)、\(G'/(\omega C')\) 不小，一阶弱损近似失效。
